% 341_GF27_GALG_FFGFT_En_ch.tex
% Johann Pascher, 1 September 2026
% Occasion: Doug Matzke, IPI mail 1 Sept. 2026

\providecommand{\BM}{\textbf{[B]}}
\providecommand{\KM}{\textbf{[K]}}
\providecommand{\SM}{\textbf{[S]}}

\chapter*{GF$(27)$ in GALG and FFGFT}
\addcontentsline{toc}{chapter}{GF$(27)$ in GALG and FFGFT}

\begin{abstract}
Doug Matzke reports in his IPI mail of 1~September~2026 that a systematic
search in GALG~G$(3)$ (6561 attempts) finds no element of order~26.
This document explains algebraically why, shows where such elements exist
in~G$(6)$, and draws the consequences for the FFGFT mass layer.

Three theorems:
\textbf{(A)}~G$(3) \cong M_2(\text{GF}(9))^2$: order~26 is algebraically
impossible in G$(3)$ --- not a search error, but a theorem.
\textbf{(B)}~G$(6) \cong M_8(\text{GF}(9))$: order-26 elements exist;
a 5-blade example can be verified directly in~GALG.
\textbf{(C)}~GF$(27)$ is present in G$(6)$ as an order structure, not as a subfield
(since $3 \nmid 8$) --- exactly the FFGFT statement from Doc.~338.

Consequence: $1/\alpha = 3700/27$ lives on the G$(3)$/GF$(9)$ level;
the lepton mass layer $(m_\mu/m_e)^2 = 43200$ needs G$(6)$ with
order-26 elements.
Verification script: \texttt{pruef\_341\_gf27\_in\_galg.py} (all assertions~\BM).
\end{abstract}

% ============================================================
\section{Background and Notation}
% ============================================================

This section provides the necessary background for readers not
immediately familiar with all the algebraic structures involved.
It can be skipped by specialists.

\subsection*{Status markers}

This document uses the standard FFGFT corpus status markers:
\begin{itemize}
  \item \textbf{[K]} — \emph{Confirmed}: algebraically or numerically fully proved;
        all assertions in the verification script passed.
  \item \textbf{[B]} — \emph{Established}: numerically secured by sampling or
        explicit construction; analytic proof present or straightforward.
  \item \textbf{[S]} — \emph{Speculative}: indication or conjecture, not yet fully proved.
\end{itemize}

\subsection*{Galois fields GF$(p^n)$}

A \emph{Galois field} GF$(p^n)$ is the unique finite field with $p^n$ elements,
where $p$ is prime (the \emph{characteristic}) and $n \geq 1$.
Its multiplicative group GF$(p^n)^* = \text{GF}(p^n) \setminus \{0\}$
is cyclic of order $p^n - 1$.
\begin{itemize}
  \item GF$(3)$: the integers $\{0,1,2\}$ modulo~3. Multiplicative group order~$2$.
  \item GF$(9) = \text{GF}(3^2)$: elements $a + bi$ with $a,b \in \{0,1,2\}$
        and $i^2 = -1 \equiv 2$. Multiplicative group order~$8$.
        This is Doug's symmetric Z3C.
  \item GF$(27) = \text{GF}(3^3)$: 27 elements, multiplicative group order~$26 = 2 \cdot 13$.
  \item GF$(81) = \text{GF}(3^4)$: 81 elements, multiplicative group order~$80 = 2^4 \cdot 5$.
\end{itemize}
The \emph{order} of an element $x$ in a group is the smallest positive $k$
with $x^k = 1$. If the group has order $m$, then every element order divides $m$
(Lagrange's theorem). This is the key fact used throughout.

\subsection*{Clifford algebra Cl$(n)$ and GALG G$(n)$}

The Clifford algebra Cl$(n)$ is generated by $n$ basis vectors $e_1,\ldots,e_n$
with the rule $e_i^2 = +1$ (hyperbits in Matzke's language) and
$e_i e_j = -e_j e_i$ for $i \neq j$.
A \emph{blade} of grade $r$ is a product $e_{i_1} e_{i_2} \cdots e_{i_r}$
with $i_1 < i_2 < \cdots < i_r$; together with the scalar~$1$ (grade~0)
they form a basis of $2^n$ elements.
\emph{GALG G$(n)$} is Cl$(n)$ over the coefficient field GF$(9) = \text{Z3C}$
(i.e.\ each blade carries a Z3C = GF$(9)$ coefficient).
As a matrix algebra:
$\text{Cl}(3) \cong M_2(\text{GF}(9))^2$ and
$\text{Cl}(6) \cong M_8(\text{GF}(9))$,
where $M_k(F)$ denotes $k \times k$ matrices over the field $F$.

\subsection*{Why GF$(27)$ and order~26?}

The lepton mass ratio $(m_\mu/m_e)^2 = 43200 = 8^2 \cdot 5^2 \cdot 27$
involves $27 = |\text{GF}(27)|$ (Doc.~338, pruef\_330/331).
The multiplicative group GF$(27)^*$ has order~$26 = 2 \cdot 13$.
An element of order~26 in any algebra generates a copy of GF$(27)^*$.
Doug's search for order-26 elements in G$(3)$ is therefore a direct
test of whether GF$(27)$ structure is present in G$(3)$.

\subsection*{Companion matrix}

Given a monic polynomial $f(x) = x^n + a_{n-1}x^{n-1} + \cdots + a_0$,
its \emph{companion matrix} is the $n \times n$ matrix with 1's on
the subdiagonal and the coefficients $-a_0, \ldots, -a_{n-1}$
in the last column. Its characteristic polynomial is $f(x)$,
so its eigenvalues are the roots of $f$ and its order in $\text{GL}_n$
equals the multiplicative order of the roots.

% ============================================================
\section{Occasion and Starting Point}
% ============================================================

Doug Matzke (IPI, 1~Sept.~2026) reports finding in GALG~G$(3)$:
GF$(9)$ in the generator $(a+ja)^8 = 1$;
GF$(81)$ in G$(3)$ with $(1+a+b+c+(a\wedge b))^{80} = 1$;
and element orders $\{2,3,4,8,16,32,40,80\}$ --- all of which divide
$80 = |\text{GF}(81)^*|$, consistent with the Cl$(3) \cong M_2(\text{GF}(9))^2$
structure where eigenvalues live in GF$(81)$.

An explicit search for elements of order~26 (which would signal GF$(27)^*$
structure):
\begin{quote}
\texttt{gasolve([a,b,c], is\_root26): Attempted 6561 with 0 found.}
\end{quote}
$6561 = 3^8$ is the number of elements in a 4-trit search space ---
exhaustive at that depth.

This is not a search error or a limitation of the tool.
It is an algebraic theorem: order~26 requires the prime factor~13,
and~13 does not appear in $80 = |\text{GF}(81)^*|$.
The following sections prove this precisely and show where
GF$(27)$ structure does appear --- in G$(6)$, not G$(3)$.

% ============================================================
\section{Theorem A: G$(3)$ Contains No GF$(27)$}
% ============================================================

\subsection{Structure of G$(3)$}

GALG~G$(n) = \text{Cl}(n)$ over the symmetric Z3C $= \text{GF}(9) = \text{GF}(3)[i]$,
where $i^2 = -1$, characteristic~3.
For odd $n = 3$:
\begin{equation}
  G(3) = \text{Cl}(3) \cong M_2(\text{GF}(9)) \oplus M_2(\text{GF}(9))
\end{equation}
Every irreducible representation is 2-dimensional over GF$(9)$.

\subsection{Eigenvalue Structure}

The eigenvalues of a $2 \times 2$ matrix over GF$(9)$ lie in the
algebraic closure --- in the smallest extension field
GF$(9^2) = \text{GF}(81)$.
\begin{equation}
  |\text{GF}(81)^*| = 80 = 2^4 \cdot 5
\end{equation}
The number~13 does not divide 80. Therefore:
\begin{itemize}
  \item No element of $M_2(\text{GF}(9))$ has an order divisible by~13.
  \item Order~26 requires $13 \mid \text{ord}(X)$ --- impossible in G$(3)$.
  \item Doug's search space $\{a, b, c\}$ spans G$(3)$ --- the result 0/6561 is algebraically forced~\BM.
\end{itemize}

\subsection{Consistency with Doug's Findings}

\begin{table}[H]
\centering
\caption{Orders in GALG G$(3)$ and their Galois origin}
\begin{tabular}{rll}
\toprule
Order & Origin & Divides \\
\midrule
2 & $|\text{GF}(3)^*|$ & 80 \\
4 & $\varphi(5) = 4$ & 80 \\
5 & min.\ prime factor $|\text{GF}(81)^*|$ & 80 \\
8 & $|\text{GF}(9)^*|$ & 80 \\
16, 20, 40, 80 & subgroups of $\text{GF}(81)^*$ & 80 \\
3 & unipotent (char.\ 3) & --- \\
26 & $|\text{GF}(27)^*|$, needs factor 13 & \textbf{not} 80 \\
\bottomrule
\end{tabular}
\end{table}

A numerical sample of 3000 random Cl$(3)$ elements confirms:
no order is divisible by~13~\BM.

% ============================================================
\section{Theorem B: G$(6)$ Contains Order-26 Elements}
% ============================================================

\subsection{Structure of G$(6)$}

\begin{equation}
  G(6) = \text{Cl}(6) \cong M_8(\text{GF}(9))
\end{equation}
Irreducible representations are 8-dimensional over GF$(9)$.

\subsection{Construction via Irreducible Polynomial}

The polynomial $f(x) = x^3 + 2x + 1$ over GF$(3)$ (equivalent to $x^3 - x + 1$)
has no root in GF$(3)$ and is therefore irreducible.
Since $\gcd(3,2) = 1$, $f$ remains irreducible over GF$(9)$;
its roots lie in GF$(27) \setminus \text{GF}(3)$.
The companion matrix
\begin{equation}
  C_f = \begin{pmatrix} 0 & 0 & 2 \\ 1 & 0 & 1 \\ 0 & 1 & 0 \end{pmatrix}
  \in M_3(\text{GF}(3))
\end{equation}
has order~26 in $\text{GL}_3(\text{GF}(3))$ ($f$ is primitive)~\BM.

\subsection{Embedding in G$(6)$ as a GALG Element}

The $8 \times 8$ element $X = C_f \oplus C_f \oplus (-I_2)$ satisfies:
\begin{itemize}
  \item $X^{26} = 1$, $X^{13} \neq 1$, $X^{2} \neq 1$ --- order exactly 26~\BM
  \item Blade decomposition over 38 Z3C blades from $\{e_1,\ldots,e_6\}$
\end{itemize}

\subsection{Minimal Example: 5 Blades, Directly Verifiable in GALG}
\label{sec:minimal_en}

A systematic search shows: order-26 elements need at least~5 blades;
with~4 blades there are none. The following 5-blade element
can be verified directly in Doug's GALG tool:
\begin{equation}
  X = -(1+i)(e_1 e_2) + (1+i)(e_1 e_2 e_3 e_6) + (1-i)(e_1 e_2 e_5 e_6)
      + (1+i)(e_1 e_4 e_5 e_6) - (e_1 e_2 e_3 e_4 e_6)
\end{equation}
GALG syntax:
\begin{quote}
\texttt{X = -(1+1j)*(e1\^{}e2) + (1+1j)*(e1\^{}e2\^{}e3\^{}e6)}\\
\texttt{  + (1-1j)*(e1\^{}e2\^{}e5\^{}e6) + (1+1j)*(e1\^{}e4\^{}e5\^{}e6)}\\
\texttt{  - (e1\^{}e2\^{}e3\^{}e4\^{}e6)}\\
\texttt{X**26 == 1  \# True}\\
\texttt{X**13 == 1  \# False}
\end{quote}
\BM\ (pruef\_341, all assertions passed)

\subsection{Frequency in G$(6)$}

\begin{table}[H]
\centering
\caption{Fraction of elements with $13 \mid \text{ord}(X)$ in G$(6)$ by blade support}
\begin{tabular}{rcc}
\toprule
Support (number of blades) & Hits / Invertible & Fraction \\
\midrule
$\leq 4$ & 0 / $\sim$450 & 0\,\% \\
5 & rare & $< 1$\,\% \\
6 & 85 / 478 & 17.8\,\% \\
10 & 190 / 511 & 37.2\,\% \\
\bottomrule
\end{tabular}
\end{table}

Order-26 elements are common in G$(6)$ --- Doug's search in G$(3)$
($\text{span}\{a,b,c\}$) could not find them for two reasons:
algebraically (Theorem~A) and because G$(3) \subsetneq G(6)$.

% ============================================================
\section{Theorem C: GF$(27)$ as Order Structure, not Subfield}
% ============================================================

\subsection{Subfield Criterion}

A unital embedding $\text{GF}(27) \hookrightarrow M_n(\text{GF}(9))$
makes $\text{GF}(9)^n$ a
$\text{GF}(27) \otimes \text{GF}(9) = \text{GF}(729)$-vector space.
Since $[\text{GF}(729):\text{GF}(9)] = 3$, we need $3 \mid n$.

\begin{table}[H]
\centering
\caption{GF$(27)$ as subfield in $M_n(\text{GF}(9))$}
\begin{tabular}{rcc}
\toprule
$n$ & $3 \mid n$ & Subfield possible? \\
\midrule
2 & no & no (G$(3)$) \\
8 & no & no (G$(6)$) \\
3 & yes & yes \\
6 & yes & yes \\
9 & yes & yes \\
\bottomrule
\end{tabular}
\end{table}

For G$(6)$ with $n = 8$: no GF$(27)$ subfield. But the
minimal polynomial of $X$ factors:
\begin{equation}
  \chi_X(x) = (x^3 + 2x + 1)^2 \cdot (x+1)^2
  \quad \in \text{GF}(3)[x]
\end{equation}
Two cubic factors --- eigenvalues in GF$(27)$ --- and one linear factor
from GF$(3)$. GF$(27)$ appears as splitting field, not as subalgebra.

\subsection{Connection to Docs.~330/338}

This is exactly the FFGFT statement from Docs.~330 and~338:
\begin{equation}
  \text{GF}(9) \not\subset \text{GF}(27) \quad (\text{since } \gcd(2,3) = 1)
\end{equation}
Both fields sit in parallel in $\text{GF}(3^6) = \text{GF}(729)$,
not in a tower. The mass-ratio identity
$(m_\mu/m_e)^2 = |\text{GF}(9)^*|^2 \cdot 5^2 \cdot |\text{GF}(27)|$
uses both fields \emph{in parallel} --- consistent with Theorem~C.

% ============================================================
\section{Consequences for FFGFT and GALG}
% ============================================================

\begin{table}[H]
\centering
\caption{GF$(27)$ structure in GALG and FFGFT}
\begin{tabular}{llll}
\toprule
Level & GALG & FFGFT & Status \\
\midrule
$\alpha$ & G$(3)$/GF$(9)$ & $1/\alpha = 3700/27$, cancels & \BM \\
Lepton masses & G$(6)$, ord.~26 needed & $(m_\mu/m_e)^2 = 43200 = 8^2\cdot5^2\cdot27$ & \KM \\
GF$(27)$ subfield & absent in G$(3)$, G$(6)$ & parallel in GF$(729)$ & \BM \\
\bottomrule
\end{tabular}
\end{table}

The fine-structure constant $1/\alpha = 3700/27$ lives on the
GF$(9)$/G$(3)$ level: $|\text{GF}(27)| = 27$ cancels in the derivation.
Doug's statement ``I'm not working on mass'' and his search result 0/6561
are both consistent with this layering:
\begin{itemize}
  \item G$(3)$: GF$(9)$, GF$(81)$, $\alpha$-structure --- Doug's working level~\BM
  \item G$(6)$: order-26 elements, lepton mass layer --- FFGFT level~\KM
\end{itemize}

% ============================================================
\section*{Appendix: For the Reader Without an Algebra Background}
\addcontentsline{toc}{section}{Appendix: For the Reader Without an Algebra Background}
% ============================================================

This appendix explains the key ideas in plain language, assuming no prior knowledge.

\subsection*{What is a finite field?}

A field is a set of numbers in which you can add, subtract, multiply and
divide (except by zero). Familiar examples are the real numbers $\mathbb{R}$
and the rational numbers $\mathbb{Q}$.
A \emph{finite} field has only finitely many elements.
You compute ``modulo'' a prime $p$: discard all multiples of $p$.
Example GF$(3)$: the numbers $\{0,1,2\}$ with $1+2 = 0$, $2+2 = 1$, etc.

\subsection*{What does GF$(p^n)$ and GF$(p^n)^*$ mean?}

GF$(p^n)$ is the unique finite field with exactly $p^n$ elements.
The \emph{star} GF$(p^n)^*$ denotes all elements \emph{except zero} ---
analogous to $\mathbb{Z}^* = \{\ldots,-1,+1,\ldots\}$ vs.\ $\mathbb{Z}$:
\begin{equation*}
  |\text{GF}(27)| = 27 \qquad |\text{GF}(27)^*| = 26 = 27-1
\end{equation*}

\subsection*{Group order and Lagrange's theorem}

GF$(p^n)^*$ is always a cyclic group of order $p^n-1$.
The \emph{order} of an element $x$ is the smallest $k$ with $x^k = 1$.
By Lagrange's theorem, every element order \emph{divides} $p^n - 1$.
Think of a clock hand that returns every 4 hours: it cannot fit on a
6-hour clock because 4 does not divide 6.

Concretely:
\begin{itemize}
  \item GF$(81)^*$: order $80 = 2^4 \cdot 5$. Prime factors: 2 and 5 only.
        Order 26 $= 2\cdot\mathbf{13}$ is impossible --- 13 does not divide 80.
  \item GF$(27)^*$: order $26 = 2 \cdot 13$. Elements of order 13 and 26 exist here.
\end{itemize}

\subsection*{Why does GF$(27)$ appear in the lepton masses?}

$(m_\mu/m_e)^2 = 43200 = 8^2 \cdot 25 \cdot 27$.
The 27 is $|\text{GF}(27)|$ (total field size including zero);
the 8 is $|\text{GF}(9)^*|$ (group order of GF$(9)$).
These arise because the FFGFT winding numbers on the T$^4$ torus
are algebraically linked to these fields.

\subsection*{Why does Doug find no GF$(27)$ in G$(3)$?}

G$(3)$ uses $2\times2$ matrices over GF$(9)$; eigenvalues live in GF$(81)$.
GF$(81)^*$ has order 80, and $13 \nmid 80$, so order 26 is impossible.
In G$(6)$ with $8\times8$ matrices, order 26 is achievable.
A 5-blade example is given in §\ref{sec:minimal_en} and directly
verifiable in GALG.

% ============================================================
\section*{Verification Script}
% ============================================================

\texttt{pruef\_341\_gf27\_in\_galg.py} --- all three theorems verified numerically
(Cl$(6)$ over GF$(9)$ as $8\times8$ matrices; exact order determination;
blade decomposition with Z3C coefficients; samples of 3000 and 600 elements).
All assertions passed~\BM.

% ============================================================
\section*{Note on AI Assistance}
% ============================================================

This document was prepared with support from Claude (Anthropic).
All algebraic computations were verified by Python scripts with exact
arithmetic over GF$(9)$ as $8\times8$ matrix representations.
The physical interpretations and conclusions are by
Johann Pascher (ORCID 0009-0000-6518-4064).

\begin{thebibliography}{9}

\bibitem{Matzke2026IPI}
D.~Matzke,
\textit{IPI mail: GF(27) in GALG},
1~September 2026 (personal correspondence).

\bibitem{Pascher330}
J.~Pascher,
\textit{Doc.~330: GF(9)-FFGFT bridge},
FFGFT corpus.

\bibitem{Pascher338}
J.~Pascher,
\textit{Doc.~338: Mass calculations and Galois comparison},
FFGFT corpus, 28~August 2026.

\bibitem{Pascher333}
J.~Pascher,
\textit{pruef\_341\_gf27\_in\_galg.py},
FFGFT corpus, 1~September 2026.

\bibitem{Matzke2022}
D.~Matzke,
\textit{G(6,$\mathbb{Z}_3\mathbb{C}$) and GALG framework},
IPI Letters 2022ff.

\end{thebibliography}
